% Physical page 2 of the concise study, which carries this sheet and nothing
% else. Every row is the generated orientation-v2 projection of
% research/chronology.toml, placed once in the generator's order, exactly as
% the reviewed Scriptural Date and Location appendix of the expansive study
% places it, with that appendix's occurrence-and-place phrases. The two rows
% that carry the votive Tract keep the rubric's label of research/scope.md
% section 8.2; the Daniel composition row points to the appendix for the later
% dating it records. The generated definitions are imported here, once,
% because the concise entrypoint's own preamble does not carry them.
% Generated. Do not edit.
% Deterministic projection of research/chronology.toml.
% schema: 5
% document: liturgy/roman-rite/1962/propers/temporal/f01-christ-the-king
% generated-by: tools/tpt proper-chronology annotations
\newcommand{\chronologyannotationclaim}[6]{#6}
\newcommand{\chronologyannotationreach}[3]{}
\newcommand{\chronologyannotationgroup}[3]{#3}
\newcommand{\chronologyannotation}[1]{%
  \ifcsname triptychchronologyannotation@#1\endcsname
    \csname triptychchronologyannotation@#1\endcsname
  \else
    \PackageError{triptych}{No chronology annotation for #1}{}%
  \fi}
% introit: Introit
\expandafter\def\csname triptychchronologyannotation@introit\endcsname{%
  \chronologyannotationgroup{historical-setting}{research-pending}{\textbf{Historical setting} -- Occasion date unresolved.}%
  \space \chronologyannotationgroup{composition}{nonuniform}{\textbf{Composition} -- No single date applies across every cited locus.}%
}
% epistle: Epistle
\expandafter\def\csname triptychchronologyannotation@epistle\endcsname{%
  \chronologyannotationgroup{composition}{disputed}{\textbf{Composition} -- disputed: \chronologyannotationreach{Col.1.12}{Col}{true}\chronologyannotationreach{Col.1.13}{Col}{true}\chronologyannotationreach{Col.1.14}{Col}{true}\chronologyannotationreach{Col.1.15}{Col}{true}\chronologyannotationreach{Col.1.16}{Col}{true}\chronologyannotationreach{Col.1.17}{Col}{true}\chronologyannotationreach{Col.1.18}{Col}{true}\chronologyannotationreach{Col.1.19}{Col}{true}\chronologyannotationreach{Col.1.20}{Col}{true}\chronologyannotationclaim{composition.epistle-to-the-colossians}{composition}{catholic-traditional-v1}{disputed}{(Philemon; Colossians; Ephesians; Philippians), 61}{A.D. 61}; \chronologyannotationreach{Col.1.12}{Col}{true}\chronologyannotationreach{Col.1.13}{Col}{true}\chronologyannotationreach{Col.1.14}{Col}{true}\chronologyannotationreach{Col.1.15}{Col}{true}\chronologyannotationreach{Col.1.16}{Col}{true}\chronologyannotationreach{Col.1.17}{Col}{true}\chronologyannotationreach{Col.1.18}{Col}{true}\chronologyannotationreach{Col.1.19}{Col}{true}\chronologyannotationreach{Col.1.20}{Col}{true}\chronologyannotationclaim{composition.epistle-to-the-colossians}{composition}{catholic-traditional-v1}{disputed}{some time between A.D. 61 and 63}{A.D. 61--63}.}%
}
% gradual: Gradual
\expandafter\def\csname triptychchronologyannotation@gradual\endcsname{%
  \chronologyannotationgroup{historical-setting}{research-pending}{\textbf{Historical setting} -- Occasion date unresolved.}%
  \space \chronologyannotationgroup{composition}{preferred}{\textbf{Composition}: \chronologyannotationreach{Ps.71.11}{Ps}{true}\chronologyannotationreach{Ps.71.8}{Ps}{true}\chronologyannotationclaim{critical.psalms.latest-composition-boundary}{composition}{catholic-critical-v1}{preferred}{before the Maccabean period, around 165 B.C.; no individual psalm can be dated securely}{Before c. 165 B.C.}}%
}
% alleluia: Alleluia
\expandafter\def\csname triptychchronologyannotation@alleluia\endcsname{%
  \chronologyannotationgroup{historical-setting}{preferred}{\textbf{Historical setting}: \chronologyannotationreach{Dan.7.14}{Dan}{true}\chronologyannotationclaim{israel.exile.seventy-years}{historical-setting}{catholic-traditional-v1}{preferred}{seventy years}{The seventy years of servitude to Babylon, ``seventy years'' (duration)}.}%
  \space \chronologyannotationgroup{composition}{preferred}{\textbf{Composition}: \chronologyannotationreach{Dan.7.14}{Dan}{true}\chronologyannotationclaim{composition.book-of-daniel}{composition}{catholic-traditional-v1}{preferred}{admits 570-536 B.C. as its date of composition}{B.C. 570--536}.}%
}
% tract: Tract
\expandafter\def\csname triptychchronologyannotation@tract\endcsname{%
  \chronologyannotationgroup{historical-setting}{preferred}{\textbf{Historical setting}: Preferred: \chronologyannotationreach{Ps.88.27}{Ps.88}{true}\chronologyannotationreach{Ps.88.28}{Ps.88}{true}\chronologyannotationreach{Ps.88.30}{Ps.88}{true}\chronologyannotationclaim{israel.exile.third-captivity}{historical-setting}{catholic-traditional-v1}{preferred}{A.M. 3416}{The third captivity of Juda: the destruction of Jerusalem under Sedecias, A.M. 3416}; alternatives: \chronologyannotationreach{Ps.88.27}{Ps.88}{true}\chronologyannotationreach{Ps.88.28}{Ps.88}{true}\chronologyannotationreach{Ps.88.30}{Ps.88}{true}\chronologyannotationclaim{israel.exile.third-captivity}{historical-setting}{catholic-traditional-v1}{alternate}{586 B.C.}{The third captivity of Juda: the destruction of Jerusalem under Sedecias, B.C. 586}, \chronologyannotationreach{Ps.88.27}{Ps.88}{true}\chronologyannotationreach{Ps.88.28}{Ps.88}{true}\chronologyannotationreach{Ps.88.30}{Ps.88}{true}\chronologyannotationclaim{israel.exile.third-captivity}{historical-setting}{catholic-traditional-v1}{alternate}{the fall of Jerusalem (587 B.C.)}{The third captivity of Juda: the destruction of Jerusalem under Sedecias, B.C. 587}, \chronologyannotationreach{Ps.88.27}{Ps.88}{true}\chronologyannotationreach{Ps.88.28}{Ps.88}{true}\chronologyannotationreach{Ps.88.30}{Ps.88}{true}\chronologyannotationclaim{israel.exile.third-captivity}{historical-setting}{catholic-traditional-v1}{alternate}{B.C. 588}{The third captivity of Juda: the destruction of Jerusalem under Sedecias, B.C. 588}.}%
  \space \chronologyannotationgroup{composition}{preferred}{\textbf{Composition}: \chronologyannotationreach{Ps.88.27}{Ps}{true}\chronologyannotationreach{Ps.88.28}{Ps}{true}\chronologyannotationreach{Ps.88.30}{Ps}{true}\chronologyannotationclaim{critical.psalms.latest-composition-boundary}{composition}{catholic-critical-v1}{preferred}{before the Maccabean period, around 165 B.C.; no individual psalm can be dated securely}{Before c. 165 B.C.}}%
}
% gospel: Gospel
\expandafter\def\csname triptychchronologyannotation@gospel\endcsname{%
  \chronologyannotationgroup{narrated-event}{research-pending}{\textbf{Event} -- No narrated-event date in the chronology corpus.}%
  \space \chronologyannotationgroup{composition}{disputed}{\textbf{Composition} -- disputed: \chronologyannotationreach{John.18.33}{John}{true}\chronologyannotationreach{John.18.34}{John}{true}\chronologyannotationreach{John.18.35}{John}{true}\chronologyannotationreach{John.18.36}{John}{true}\chronologyannotationreach{John.18.37}{John}{true}\chronologyannotationclaim{composition.gospel-of-john}{composition}{catholic-traditional-v1}{disputed}{from the year 90 to 100 (approximately)}{c. A.D. 90--100}; \chronologyannotationreach{John.18.33}{John}{true}\chronologyannotationreach{John.18.34}{John}{true}\chronologyannotationreach{John.18.35}{John}{true}\chronologyannotationreach{John.18.36}{John}{true}\chronologyannotationreach{John.18.37}{John}{true}\chronologyannotationclaim{composition.gospel-of-john}{composition}{catholic-traditional-v1}{disputed}{96 or one of the succeeding years}{A.D. 96--100}.}%
}
% offertory: Offertory
\expandafter\def\csname triptychchronologyannotation@offertory\endcsname{%
  \chronologyannotationgroup{historical-setting}{research-pending}{\textbf{Historical setting} -- Occasion date unresolved.}%
  \space \chronologyannotationgroup{composition}{preferred}{\textbf{Composition}: \chronologyannotationreach{Ps.2.8}{Ps}{true}\chronologyannotationclaim{critical.psalms.latest-composition-boundary}{composition}{catholic-critical-v1}{preferred}{before the Maccabean period, around 165 B.C.; no individual psalm can be dated securely}{Before c. 165 B.C.}}%
}
% communion: Communion
\expandafter\def\csname triptychchronologyannotation@communion\endcsname{%
  \chronologyannotationgroup{historical-setting}{research-pending}{\textbf{Historical setting} -- Occasion date unresolved.}%
  \space \chronologyannotationgroup{composition}{preferred}{\textbf{Composition}: \chronologyannotationreach{Ps.28.10}{Ps}{true}\chronologyannotationreach{Ps.28.11}{Ps}{true}\chronologyannotationclaim{critical.psalms.latest-composition-boundary}{composition}{catholic-critical-v1}{preferred}{before the Maccabean period, around 165 B.C.; no individual psalm can be dated securely}{Before c. 165 B.C.}}%
}
% presentation: orientation-v2
\newcommand{\chronologyorientationrow}[2]{%
  \ifcsname triptychchronologyorientation@#1\endcsname
    \csname triptychchronologyorientation@#1\endcsname{#2}%
  \else
    \PackageError{triptych}{No chronology orientation row for #1}{}%
  \fi}
\expandafter\def\csname triptychchronologyorientation@row-6c43dd5e0531017f\endcsname#1{%
Introit: Ps 71:1 (Vulgate numbering; Hebrew 72); Gradual: Psalms 71:8,11 (Vulgate numbering; Hebrew 72); Offertory: Psalms 2:8 (Vulgate and Hebrew numbering agree); Communion: Psalms 28:10--11 (Vulgate numbering; Hebrew 29) & Historical setting. #1 & Occasion date unresolved\\[.25em]%
}
\expandafter\def\csname triptychchronologyorientation@row-7fbe70097919820b\endcsname#1{%
Introit: Ps 71:1 (Vulgate numbering; Hebrew 72); Gradual: Psalms 71:8,11 (Vulgate numbering; Hebrew 72); Tract: Psalms 88:27--28,30 (Vulgate numbering; Hebrew 89); Offertory: Psalms 2:8 (Vulgate and Hebrew numbering agree); Communion: Psalms 28:10--11 (Vulgate numbering; Hebrew 29) & Composition: Psalms. #1 & \chronologyannotationreach{Ps.2.8}{Ps}{true}\chronologyannotationreach{Ps.28.10}{Ps}{true}\chronologyannotationreach{Ps.28.11}{Ps}{true}\chronologyannotationreach{Ps.71.1}{Ps}{true}\chronologyannotationreach{Ps.71.11}{Ps}{true}\chronologyannotationreach{Ps.71.8}{Ps}{true}\chronologyannotationreach{Ps.88.27}{Ps}{true}\chronologyannotationreach{Ps.88.28}{Ps}{true}\chronologyannotationreach{Ps.88.30}{Ps}{true}\chronologyannotationclaim{critical.psalms.latest-composition-boundary}{composition}{catholic-critical-v1}{preferred}{before the Maccabean period, around 165 B.C.; no individual psalm can be dated securely}{Before c. 165 B.C.}\\[.25em]%
}
\expandafter\def\csname triptychchronologyorientation@row-3bd03395c37c5ada\endcsname#1{%
Tract: Psalms 88:27--28,30 (Vulgate numbering; Hebrew 89) & Historical setting: Jerusalem's destruction under Sedecias. #1 & Preferred: \chronologyannotationreach{Ps.88.27}{Ps.88}{true}\chronologyannotationreach{Ps.88.28}{Ps.88}{true}\chronologyannotationreach{Ps.88.30}{Ps.88}{true}\chronologyannotationclaim{israel.exile.third-captivity}{historical-setting}{catholic-traditional-v1}{preferred}{A.M. 3416}{A.M. 3416 (Haydock reporting Ussher)}; alternatives: \chronologyannotationreach{Ps.88.27}{Ps.88}{true}\chronologyannotationreach{Ps.88.28}{Ps.88}{true}\chronologyannotationreach{Ps.88.30}{Ps.88}{true}\chronologyannotationclaim{israel.exile.third-captivity}{historical-setting}{catholic-traditional-v1}{alternate}{586 B.C.}{B.C. 586 (Reid and Meistermann)}, \chronologyannotationreach{Ps.88.27}{Ps.88}{true}\chronologyannotationreach{Ps.88.28}{Ps.88}{true}\chronologyannotationreach{Ps.88.30}{Ps.88}{true}\chronologyannotationclaim{israel.exile.third-captivity}{historical-setting}{catholic-traditional-v1}{alternate}{the fall of Jerusalem (587 B.C.)}{B.C. 587 (Schets)}, \chronologyannotationreach{Ps.88.27}{Ps.88}{true}\chronologyannotationreach{Ps.88.28}{Ps.88}{true}\chronologyannotationreach{Ps.88.30}{Ps.88}{true}\chronologyannotationclaim{israel.exile.third-captivity}{historical-setting}{catholic-traditional-v1}{alternate}{B.C. 588}{B.C. 588 (Petavius via Sloet; transcription only)}\\[.25em]%
}
\expandafter\def\csname triptychchronologyorientation@row-94f74a8fe6bde9b6\endcsname#1{%
Alleluia: Daniel 7:14 & Historical setting: The seventy years of servitude to Babylon. #1 & \chronologyannotationreach{Dan.7.14}{Dan}{true}\chronologyannotationclaim{israel.exile.seventy-years}{historical-setting}{catholic-traditional-v1}{preferred}{seventy years}{``seventy years'' (duration)}\\[.25em]%
}
\expandafter\def\csname triptychchronologyorientation@row-98256791fb1a09b1\endcsname#1{%
Alleluia: Daniel 7:14 & Composition: The Book of Daniel. #1 & \chronologyannotationreach{Dan.7.14}{Dan}{true}\chronologyannotationclaim{composition.book-of-daniel}{composition}{catholic-traditional-v1}{preferred}{admits 570-536 B.C. as its date of composition}{B.C. 570--536}\\[.25em]%
}
\expandafter\def\csname triptychchronologyorientation@row-bb05ffe36adb121c\endcsname#1{%
Gospel: John 18:33--37 & Event. #1 & Date unresolved\\[.25em]%
}
\expandafter\def\csname triptychchronologyorientation@row-b5624d4c980f37b6\endcsname#1{%
Gospel: John 18:33--37 & Composition: The Gospel of St John. #1 & Disputed: \chronologyannotationreach{John.18.33}{John}{true}\chronologyannotationreach{John.18.34}{John}{true}\chronologyannotationreach{John.18.35}{John}{true}\chronologyannotationreach{John.18.36}{John}{true}\chronologyannotationreach{John.18.37}{John}{true}\chronologyannotationclaim{composition.gospel-of-john}{composition}{catholic-traditional-v1}{disputed}{from the year 90 to 100 (approximately)}{c. A.D. 90--100}; \chronologyannotationreach{John.18.33}{John}{true}\chronologyannotationreach{John.18.34}{John}{true}\chronologyannotationreach{John.18.35}{John}{true}\chronologyannotationreach{John.18.36}{John}{true}\chronologyannotationreach{John.18.37}{John}{true}\chronologyannotationclaim{composition.gospel-of-john}{composition}{catholic-traditional-v1}{disputed}{96 or one of the succeeding years}{A.D. 96--100}\\[.25em]%
}
\expandafter\def\csname triptychchronologyorientation@row-e679cfff41e9473d\endcsname#1{%
Epistle: Colossians 1:12--20 & Composition: The Epistle to the Colossians. #1 & Disputed: \chronologyannotationreach{Col.1.12}{Col}{true}\chronologyannotationreach{Col.1.13}{Col}{true}\chronologyannotationreach{Col.1.14}{Col}{true}\chronologyannotationreach{Col.1.15}{Col}{true}\chronologyannotationreach{Col.1.16}{Col}{true}\chronologyannotationreach{Col.1.17}{Col}{true}\chronologyannotationreach{Col.1.18}{Col}{true}\chronologyannotationreach{Col.1.19}{Col}{true}\chronologyannotationreach{Col.1.20}{Col}{true}\chronologyannotationclaim{composition.epistle-to-the-colossians}{composition}{catholic-traditional-v1}{disputed}{(Philemon; Colossians; Ephesians; Philippians), 61}{A.D. 61}; \chronologyannotationreach{Col.1.12}{Col}{true}\chronologyannotationreach{Col.1.13}{Col}{true}\chronologyannotationreach{Col.1.14}{Col}{true}\chronologyannotationreach{Col.1.15}{Col}{true}\chronologyannotationreach{Col.1.16}{Col}{true}\chronologyannotationreach{Col.1.17}{Col}{true}\chronologyannotationreach{Col.1.18}{Col}{true}\chronologyannotationreach{Col.1.19}{Col}{true}\chronologyannotationreach{Col.1.20}{Col}{true}\chronologyannotationclaim{composition.epistle-to-the-colossians}{composition}{catholic-traditional-v1}{disputed}{some time between A.D. 61 and 63}{A.D. 61--63}\\[.25em]%
}
\expandafter\def\csname triptychchronologyorientation@row-1dc5d7a61ec35b95\endcsname#1{%
Introit: Apoc 1:6, Apoc 5:12 & Composition: The Apocalypse of St John. #1 & Disputed: \chronologyannotationreach{Apoc.1.6}{Apoc}{true}\chronologyannotationreach{Apoc.5.12}{Apoc}{true}\chronologyannotationclaim{composition.apocalypse}{composition}{catholic-traditional-v1}{disputed}{the reign of Domitian ... A.D. 81-96}{A.D. 81--96}; \chronologyannotationreach{Apoc.1.6}{Apoc}{true}\chronologyannotationreach{Apoc.5.12}{Apoc}{true}\chronologyannotationclaim{composition.apocalypse}{composition}{catholic-traditional-v1}{disputed}{from the year 90 to 100 (approximately)}{c. A.D. 90--100}\\[.25em]%
}


\runninghead{Scriptural Date and Location}
\section{Scriptural Date and Location}

\zlabel{triptych:concise:chronology:start}
\begin{chronologyorientationtable}
\chronologyorientationrow{row-6c43dd5e0531017f}{Psalms of a king and of the Lord as King. Ps 71 is titled ``on Solomon'' and closes with the praises of David (71:20); Ps 2, where the king is set on Sion, is ascribed to David in Acts 4:25; Ps 28, which ends with the Lord sitting as King for ever, is titled ``for David, at the finishing of the tabernacle'' (28:1). No occasion is dated.}
\addlinespace
\chronologyorientationrow{row-7fbe70097919820b}{The psalms of this Mass, under the titles above. Ps 88, titled ``for Ethan the Ezrahite'' (88:1), is the psalm of the Tract, printed for votive Masses after Septuagesima, never said on this feast.}
\addlinespace
\chronologyorientationrow{row-3bd03395c37c5ada}{The Tract, printed for votive Masses after Septuagesima, never said on this feast. The setting dated is the psalm's lament that God has rejected his anointed and cast his throne down to the ground (Ps 88:39--45).}
\midrule
\chronologyorientationrow{row-94f74a8fe6bde9b6}{Daniel's night vision ``in the first year of Baltasar, king of Babylon'' (Dan 7:1), among the exiles of Juda in Babylon.}
\addlinespace
\chronologyorientationrow{row-98256791fb1a09b1}{The book that presents Daniel, a youth of Juda carried to Babylon (Dan 1:1--6). The same verse opens the paschal Alleluia, printed for Paschaltide, never said on this feast; its second verse, Apoc 19:16, has no row. A later dating is set out in the expansive study's Scriptural Date and Location.}
\midrule
\chronologyorientationrow{row-bb05ffe36adb121c}{Jesus questioned by Pilate inside the governor's hall, the praetorium, at Jerusalem, on the morning before the Pasch (Jn 18:28, 33).}
\addlinespace
\chronologyorientationrow{row-b5624d4c980f37b6}{John writes his Gospel, by early tradition at Ephesus (Irenaeus, as Fonck reports him).}
\midrule
\chronologyorientationrow{row-e679cfff41e9473d}{Paul writes in bonds (Col 4:3, 18) to the saints at Colossa (1:2); the common view places the letters of the captivity at Rome.}
\midrule
\chronologyorientationrow{row-1dc5d7a61ec35b95}{John writes the visions he saw ``in the island which is called Patmos'' (Apoc 1:9).}
\end{chronologyorientationtable}
\zlabel{triptych:concise:chronology:end}
